\chapter*{附录一 《圣祖仁皇帝实录》中记载废黜太子缘由}\addcontentsline{toc}{chapter}{附录一 《圣祖仁皇帝实录》中记载废黜太子缘由}\markboth{附录一 《圣祖仁皇帝实录》中记载废黜太子缘由}{附录一 《圣祖仁皇帝实录》中记载废黜太子缘由}
% 原书第375页

（康熙四十七年九月丁丑）上召诸王、大臣、侍卫、文武官员等，齐集行宫前。命皇太子允礽跪。上垂涕，谕曰：


% 原书第375页

“朕承太祖、太宗、世祖弘业四十八年于兹，兢兢业业，轸恤臣工，惠养百姓，惟以治安天下为务。


% 原书第375页

今观允礽，不法祖德，不遵朕训，惟肆恶虐众，暴戾淫乱，难出诸口。朕包容二十年矣，乃其恶愈张。僇辱在廷诸王、贝勒、大臣、官员，专擅威权。鸠聚党与，窥伺朕躬，起居动作，无不探听。朕思国惟一主，允礽何得将诸王、贝勒、大臣、官员任意凌虐、恣行捶挞耶？如平郡王讷尔素、贝勒海善、公普奇，俱被伊殴打。大臣官员以至兵丁，鲜不遭其荼毒。朕深悉此情，因诸臣有言及伊之行事者，伊即仇视其人，横加鞭笞，故朕未将伊之行事，一询及于诸臣。


% 原书第375–376页

朕巡幸陕西、江南、浙江等处，或驻庐舍，或御舟航，未尝跬步妄出，未尝一事扰民。乃允礽同伊属下人等，恣行乖戾，无所不至，令朕赧于启齿。又遣使邀截外藩入贡之人，将进御马匹任意攘取，以至蒙古俱不心服。种种恶端，不可枚举。朕尚冀其悔过自新，故隐忍优容至于今日。


% 原书第376页

又，朕知允礽赋性奢侈，着伊乳母之父凌普为内务府总管，俾伊便于取用。孰意凌普更为贪婪，致使包衣下人无不怨恨。朕自允礽幼时，谆谆教训，凡所用物，皆系庶民脂膏，应从节俭。乃不遵朕言，穷奢极欲，逞其凶恶，今更滋甚，有将朕诸子不遗噍类之势。


% 原书第376页

十八阿哥患病，众皆以朕年高，无不为朕忧虑。伊系亲兄，毫无友爱之意，因朕加责让，伊反忿然发怒。更可异者，伊每夜逼近布城，裂缝向内窃视。从前索额图助伊潜谋大事，朕悉知其情，将索额图处死。今允礽欲为索额图复仇，结成党羽，令朕未卜今日被鸩，明日遇害，昼夜戒慎不宁。似此之人，岂可付以祖宗弘业？且允礽生而克母，此等之人，古称不孝。


% 原书第376页

朕即位以来，诸事节俭，身御敝褥，足用布袜。允礽所用，一切远过于朕，伊犹以为不足。恣取国帑，干预政事，必致败坏我国家，戕贼我万民而后已。若以此不孝不仁之人为君，其如祖业何？”


% 原书第376页

谕毕。上复痛哭仆地，诸大臣扶起。上又谕曰：


% 原书第377页

“太祖、太宗、世祖之缔造勤劳，与朕治平之天下，断不可以付此人。俟回京昭告于天地、宗庙，将允礽废斥。


% 原书第377页

朕前命直郡王允禔，善护朕躬，并无欲立允禔为皇太子之意，允禔秉性躁急愚顽，岂可立为皇太子。其允礽党羽，凡系畏威附合者皆从宽不究外，将索额图之子格尔芬、阿尔吉善暨二格、苏尔特、哈什太、萨尔邦阿，俱立行正法。杜默臣、阿进泰、苏赫陈、倪雅汉，着充发盛京。此事关系天下万民，甚属紧要。乘朕身体康健，定此大事。着将允礽即行拘执。


% 原书第377页

尔诸王、大臣、官员、兵民等，以允礽所行之事为虚为实，可各秉公陈奏。”


% 原书第377页

众皆叩首流涕，奏曰：“皇上所见、至圣至明。谕旨所言皇太子诸事，一一皆确实。臣等实无异辞可以陈奏。”

